% Experiments - ICML style

\paragraph{Data: WorldTrace segments and road graphs.}
We evaluate on WorldTrace trip segments for Detroit and Columbus.
Each segment is resampled to a fixed-length trajectory ($F=256$) and retains a Unix epoch departure time $t_0$.
We build a road graph per city from OpenStreetMap and map-match each segment to a node sequence using nearest-node snapping plus bridging for non-adjacent jumps.
For multi-city experiments, we form a disjoint union of city graphs by node-index offsetting and record the city ID per route.
\Cref{tab:data-stats} summarizes the dataset statistics.

\begin{table}[t]
  \caption{Dataset statistics for WorldTrace segments.}
  \label{tab:data-stats}
  \vskip 0.1in
  \centering
  \small
  \begin{tabular}{lcccc}
    \toprule
    City & Routes & Graph Nodes & Graph Edges & Median Path Len \\
    \midrule
    Detroit & 2,279 & 10,804 & 26,995 & 35 \\
    Columbus & 5,223 & 8,747 & --- & 22 \\
    Combined & 7,502 & 19,551 & --- & 26 \\
    \bottomrule
  \end{tabular}
  \vskip -0.1in
\end{table}

\paragraph{Diagnostic Gate 1: Candidate-set coverage.}
We diagnose the ``classify among $K$ candidates'' strategy by generating $K$ heuristic candidates (Yen-style $K$-shortest paths) per OD pair and computing best-of-$K$ edge Jaccard match to the ground-truth graph path.
We also report the distribution of ground-truth point-to-road distances to verify that low Jaccard is not caused by systematic GT/map misalignment.

\paragraph{Diagnostic Gate 2: Semantic informativeness.}
To test whether context can explain corridor variability, we group routes into coarse OD buckets and cluster the ground-truth graph paths within each bucket by hierarchical clustering on edge-Jaccard distance.
For OD buckets with at least two sufficiently large clusters, we build a binary classification task (top-2 clusters) and fit a lightweight logistic regression using time features (hour-of-day, day-of-week) and aggregated tier features.
We report AUC as a Go/No-Go signal for ``context-conditioned diversity.''

\paragraph{Baselines and methods.}
We evaluate three decision mechanisms for generating corridors on the road graph:
\begin{itemize}
  \item \textbf{$K$-shortest candidates:} Yen's algorithm; evaluated via best-of-$K$ Jaccard (diagnostic baseline).
  \item \textbf{Node-level AR:} Predict the next node among neighbors conditioned on time and tier; diversity via temperature sampling.
  \item \textbf{Waypoint AR + A* (ours):} Sample a short waypoint sequence in 3--5 steps, then connect waypoints by A* to produce a valid graph path.
\end{itemize}

\paragraph{Metrics.}
Our primary corridor-level metrics are:
\begin{itemize}
  \item \textbf{Success rate:} Fraction of sampled routes that reach $d$ within a maximum step budget.
  \item \textbf{Best-of-$K$ corridor match:} $J_{\mathrm{best}}$ computed on directed edge sets.
\end{itemize}
All reported numbers use $K=20$ samples per route.
We emphasize that $J_{\mathrm{best}}$ is a single-trajectory match metric (one GT route vs $K$ generated routes) and should not be misread as evidence of corridor diversity.
For corridor-level multi-modality analysis, we explicitly use OD buckets with multiple ground-truth instances.
